\chapter{扩散生成的计算动力学：噪声过程、条件引导与结构保持}
\label{chp:diffusion_thermodynamics}

扩散模型把复杂数据分布的采样问题改写为一条可学习的随机演化路径：训练时逐步扰动样本，使中间分布趋向便于采样的参考分布；生成时再估计每个噪声尺度上的反向更新，把随机初态变成满足数据统计和条件约束的输出。这一机制足以解释其工程效力，无须把像素噪声说成焦耳热、把得分函数说成引力，或由生成结果逼真推出模型与物理世界整体同构。本章以计算动力学为主线，分别重建前向扰动、反向采样、文本条件、时空结构失效和结构—内容联合生成，明确概率对象、离散步长、训练假设、可测误差与应用边界，并把物理语言限制在确有量纲和守恒模型的领域环节。

\section{前向扰动的概率动力学：对象、尺度与可逆信息边界}
\label{sec:forward_process_heat_death}

我们先规定扩散过程处理的对象。设数据变量 $x_0\in\mathcal{X}\subseteq\mathbb{R}^d$，其中 $d$ 是编码维数；图像可以把像素与通道展平为该向量，潜空间模型则令 $x_0$ 为编码器输出。坐标本身通常是归一化数值，没有焦耳、开尔文或米等物理单位；只有当某个坐标明确表示深度、速度或其他领域量时，才另外携带相应单位。令离散扩散步为 $k\in\{0,\ldots,K\}$，它是算法索引而非现实时间；噪声计划 $\beta_k\in(0,1)$ 无量纲，记 $\alpha_k=1-\beta_k$、$\bar\alpha_k=\prod_{j=1}^{k}\alpha_j$。经典高斯前向核定义为
\begin{equation}
q(x_k\mid x_{k-1})=\mathcal{N}\!\left(x_k;\sqrt{\alpha_k}x_{k-1},\beta_k I_d\right),
\end{equation}
并由高斯闭包得到
\begin{equation}
q(x_k\mid x_0)=\mathcal{N}\!\left(x_k;\sqrt{\bar\alpha_k}x_0,(1-\bar\alpha_k)I_d\right),\qquad
x_k=\sqrt{\bar\alpha_k}x_0+\sqrt{1-\bar\alpha_k}\,\epsilon,
\end{equation}
其中 $\epsilon\sim\mathcal{N}(0,I_d)$。这里第一式是模型定义，第二式是该定义在独立高斯增量假设下的数学结果。它说明任意训练步都可直接采样，而不必实际执行前面所有步骤；也说明当 $\bar\alpha_K$ 足够小时，终端分布才近似标准高斯。若数据没有适当中心化、协方差强烈失配、$K$ 太小或噪声计划不合适，所谓“终端纯噪声”只是一项近似，不能被写成无条件定理。

连续表述需要再声明极限。令算法时间 $t\in[0,T]$，单位可取归一化步，状态过程 $X_t\in\mathbb{R}^d$ 满足 Itô 随机微分方程
\begin{equation}
dX_t=f(X_t,t)\,dt+g(t)\,dW_t,
\end{equation}
其中漂移 $f$ 的单位为“状态/算法时间”，扩散系数 $g$ 的单位为“状态/算法时间平方根”，$W_t$ 是 $d$ 维 Wiener 过程。离散算法只有在步长趋小、系数满足常见的可积性与正则条件时，才可由该连续模型近似；实际采样器还要声明 Euler--Maruyama、预测—校正或高阶求解器、步长表和截断误差。把 $t$ 称为物理时间会混淆三个不同对象：图像或视频内部的场景时间、训练样本的记录时间，以及噪声调度的算法时间。尤其对视频生成，场景帧索引属于数据坐标，而扩散步属于求解坐标，两者不可互换。

旧稿把前向加噪称作“热寂”、把图像称作低熵拓扑孤立子，并进一步断言流形度量执行里奇流。我们保留其中“逐步抹去可辨识信息”的直观，但更换其可检验对象。对两个初始样本 $x_0,x'_0$，条件均值之差随步数缩为 $\sqrt{\bar\alpha_k}(x_0-x'_0)$；在选定距离、判别器或互信息估计器下，可以测量样本身份的可区分度怎样下降。若研究微分熵，还必须说明坐标变换和参考测度，因为连续变量的微分熵不是坐标不变量，也不等于热力学熵。加噪需要计算资源但公式中的方差不是焦耳能耗；两者只能通过硬件测量另建关联模型。类似地，卷积或注意力网络诱导的表示几何可以用局部敏感度、雅可比谱或任务距离研究，却不会仅由高斯核自动得到黎曼曲率，更不会由像素平滑推出度量满足 $\partial_tg_{\mu\nu}=-2R_{\mu\nu}$。

前向过程的工程目的不是“先毁灭万物”，而是设计一族难度连续、监督可构造的中间分布。训练可从数据集抽取 $x_0$，再抽取 $k$ 与 $\epsilon$，使网络预测噪声、干净样本、速度参数化或得分；这些目标在权重和参数化正确时可以互相换算，但有限模型、有限数据和数值精度会使其优化行为不同。以“骑马的宇航员”为例，前向核只把一幅已经存在的训练图像扰动为不同信噪比的版本，它并没有自动标注“宇航员”“马”、主体—坐骑关系或遮挡顺序。此类结构必须来自数据、条件编码、网络归纳偏置和损失设计。反例也很直接：如果训练数据几乎没有正确骑乘姿态，即使终端分布完美接近高斯，反向模型仍可能把腿、缰绳和马身绑定错误；因此参考分布易采样不等于目标组合已学会。

在 HSF--HD 的层级中，我们把前向扩散定位为一般状态动力学中的一种人为扰动协议，而不是智能系统普遍必有的“热力学毁灭阶段”。其状态记录至少包含 $(x_k,k,\nu,\sigma)$：$\nu$ 指定编码器、数据预处理和噪声计划的版本，$\sigma$ 记录随机种子或可复现实验所需的噪声来源。对局部结构网络，还应保留对象身份、邻接或区域掩码等结构变量；对全局分布表示，则记录连续特征及其不确定性。二者可以在训练时共同受扰动，也可以采用不同噪声核，但必须通过任务相关身份、绑定和读出接口重新联合。只有把这些对象和版本写清楚，我们才可能比较噪声计划、诊断信息在哪个尺度丢失，并为下一节反向过程的可识别条件建立真正的计算基线。
